% Initial plan using the qproposal class
%
% Build: pdflatex plan && bibtex plan && pdflatex plan && pdflatex plan
% Or with latexmk: latexmk -pdf plan

\documentclass[sans,oneside,single,parskip,raggedright,numeric]{qproposal}

% Tables
\usepackage{booktabs}
\usepackage{longtable}
\usepackage{array}
\usepackage{enumitem}

% ------------------------------------------------------------------------------
% Metadata
% ------------------------------------------------------------------------------

\title{Initial Plan: Utilising Generative AI to Create Vulnerabilities in Machines for Education}
\author{Haris Khan (23040365)}
\supervisor{Amir Javed}
\date{February 2026}

% ------------------------------------------------------------------------------
% Document body
% ------------------------------------------------------------------------------

\begin{document}

\maketitle

\begin{abstract}
  This document outlines the initial plan for an automated information security training
  platform that uses a large language model (LLM) to generate vulnerabilities in ephemeral,
  sandboxed environments. It describes the problem context and motivation, states the aims
  and objectives, discusses feasibility constraints, and presents a week-by-week work plan
  with milestones and deliverables.
\end{abstract}

% ------------------------------------------------------------------------------
\section{Project Description}
% ------------------------------------------------------------------------------

\subsection*{Context and motivation}

I have a significnat interest in information security, and the diversity of attack vectors.
To improve skills in this realm, hands-on experience is the best way to acquire new knowledge.
Creating a simulated environment, where users can safely practise these skills was something
that interested me as it requires careful isolation of sandboxed environments, automatically
deploying containers; and providing a safe place to hone your skills. In addition to this,
the emergence of LLMs provides an opportunity to modernize these environments to allow
infinite customisation and guidance.

\subsection*{Outline of the problem}

The problem with existing educational tools such as HackTheBox~\citep{hackthebox} and
TryHackMe~\citep{tryhackme} is the static nature of the vulnerabilities and with no direct
assistance in solving them, it makes it difficult to make progress once stuck. Therefore, a
solution is where the vulnerabilities are more dynamic, able to be tailored by the user, and
allow a personalised level of guidance rather than a bullet-point guide. Using generative AI,
you can customise the level of difficulty, and also use it as an assistant when solving the
challenge to help probe closer to the solution.

\subsection*{The specific version of it that I intend to address}

A version of the problem where the sandboxed environment is able to be tailored to the
whatever the user desires, whether it be a simple nginx misconfiguration, or a logic error
inside an web application. Even if they select the same type of vulnerability, due to the
polymorphic nature of the LLM, the exploitation method will be different. Furthermore, to
address the guidance and assistance access, it will leverage the LLMs that will be creating
the vulnerabilities in order to provide a chatbot for the specific room, with context to it,
that can help guide the user to the solution, without explicitly providing it.

\subsection*{How I intend to approach the solution}

To simplify the project, the platform will rely on SSH as an access point. Completion will
be via a `flag'~\citep{ctf-wikipedia} system, similar to those used in existing CTFs.

\subsection*{Background material}

There are many existing CTFs that are used such as picoCTF~\citep{picoctf},
VulnHub~\citep{vulnhub}, however these providers do not use generative AI to create the
sandboxes. Some provide an LLM to give assistance, but not as the creator of the
vulnerability, limiting the potential impact they have.

\subsection*{Why the problem is a sufficient challenge for the project module}

My project will cover virtualization orchestration with type-1 hypervisors, requiring a
state machine to manage the lifecycle of these machines, and resource handling with clean-up
of them. There must be strict network segmentation, and escape prevention to ensure the user
is unable to access any unintended resources. In addition to this, I will have to experiment
and develop quite complex system prompts to enable LLMs to generate vulnerabilites, which
with current guardrails may prove difficult.

\newpage

% ------------------------------------------------------------------------------
\section{Aims and Objectives}
% ------------------------------------------------------------------------------

The overall aim is to implement an automated information security training platform that
utilises an LLM to configure vulnerabilites in ephemeral environments.

\begin{itemize}
  \item \textbf{Objective 1 --- Orchestration controller.}
    Design a controller that will orchestrate and configure the containers. This will include
    having a `master' template, with dynamic configuration based on the vulnerability type,
    handling of network interfaces, and injection of users' SSH keys to handle access. Risks
    include a misconfigured environment and the LLM exceeding its scope beyond the container.
    These will be mitgated with zero-trust access~\citep{ztna-microsoft} and kill switches.

  \item \textbf{Objective 2 --- Generative vulnerability engine.}
    Implement the vulnerability engine using a local or cloud LLM (depending on resource
    constraints) to generate scripts and applications that intentionally contain a flaw,
    then inject them into the container. Risks include hallucinated or non-functional
    vulnerabilities, and model guardrails preventing generation of security flaws.

  \item \textbf{Objective 3 --- User interface and flag system.}
    Develop a web dashboard to manage machines, user accounts, and SSH connections, together
    with a flag-based completion system. Risks include SSH key management, session-state
    handling, and users struggling to reach web applications inside containers via SSH
    tunnelling.

  \item \textbf{Objective 4 --- LLM tutor.}
    Integrate a context-aware LLM tutor into the web UI as a chatbot. Because the same LLM
    session creates the vulnerability, it will have the full context of the challenge. Risks
    include the tutor leaking the flag (mitigated by system-level prompts) and handling
    large contexts; models with large context windows such as Gemini 3
    Pro~\citep{gemini-3-pro} (1 million input tokens) may prove advantageous here.
\end{itemize}

\newpage

% ------------------------------------------------------------------------------
\section{Feasibility}
% ------------------------------------------------------------------------------

The main restraint for this project is hardware capable of running a type-1 hypervisor with
multiple containers and, ideally, a local LLM. I have a personal machine I can use for the
hypervisor; for the LLM I will require cloud credits or a machine capable of running local
models. Some larger local models such as Kimi K2.5 require large amounts of RAM
($>$240\,GB~\citep{kimi-k25}), but smaller models like gpt-oss:20B can run on less powerful
hardware. There will be some loss in reasoning capability that would have to be considered,
such as struggling to create certain vulnerabilities.

Regarding ethical considerations, I will have to ensure that the models used are permitted to
create vulnerabilites for educational purposes. Protocols must be in place to prevent
vulnerable containers from being accessible externally, such as a time-based kill switch.
Regarding licensing, many open-weight models carry community licences such as Llama
3~\citep{llama3-license} which do not pose significant risk. Open-source libraries such as
Proxmox~\citep{proxmox} will be used throughout.

\newpage 

% ------------------------------------------------------------------------------
\section{Work Plan}
% ------------------------------------------------------------------------------

% Column type: paragraph with top alignment and small font
\newcolumntype{P}[1]{>{\small\RaggedRight\arraybackslash}p{#1}}
\newcolumntype{C}[1]{>{\small\centering\arraybackslash}p{#1}}

% Compact itemize for use inside table cells
\newenvironment{cellist}{%
  \begin{itemize}[nosep,leftmargin=1em,topsep=0pt,partopsep=0pt]
}{%
  \end{itemize}
}

\begin{longtable}{C{1.3cm} P{6.2cm} P{6.2cm}}
  \toprule
  \textbf{Weeks} & \textbf{Tasks} & \textbf{Milestones and deliverables} \\
  \midrule
  \endfirsthead
  \toprule
  \textbf{Weeks} & \textbf{Tasks} & \textbf{Milestones and deliverables} \\
  \midrule
  \endhead
  \bottomrule
  \endfoot

  1--3 &
  \begin{cellist}
    \item Create the initial plan and find any constraints
    \item Setup the hypervisor and the physical machine where everything will be virtualized on
    \item Create the `master' template that will be used across generated machines
    \item Configure network access and security
    \item Ensure remote access and development is possible from my local machine
    \item Verify resource allocation for the LLM
  \end{cellist}
  &
  Initial report complete and limitations identified. Proxmox reachable from local machines; able to spin up machines from the master template. Machines isolated from each other and not easily escapable.
  \\[4pt]

  3--4 &
  \begin{cellist}
    \item Develop the `orchestrator' of the virtual machines (script or Ansible playbooks with cloud-init)
    \item The orchestrator should allow access through SSH
  \end{cellist}
  &
  Orchestrator spins up machines on demand, clones the master template, and injects a specified SSH key. Accessible from an external machine.
  \\[4pt]

  5--6 &
  \begin{cellist}
    \item Test different LLM models for vulnerability generation
    \item Develop specific prompts to create different vulnerability types, and a system-wide prompt for guardrails
    \item Implement the injection pipeline so the LLM can alter containers
  \end{cellist}
  &
  Chosen model capable of generating different vulnerability types from prompts. Injection pipeline functional; vulnerability exists on the container after provisioning.
  \\[4pt]

  7--8 &
  \begin{cellist}
    \item Develop the dashboard (account creation, SSH key entry, vulnerability selection, machine provisioning)
    \item Enable SSH access to containers using functionality from weeks 3--4
    \item Implement the LLM tutor chatbot
    \item Allow the user to submit and verify the flag
  \end{cellist}
  &
  Full user login flow working. Container provisioned with selected vulnerability and accessible via SSH. Tutor chatbot answers questions with prior vulnerability context. Flag submission and verification functional.
  \\[4pt]

  8--9 &
  \begin{cellist}
    \item Test and verify site functionality, specifically that flags are achievable and vulnerabilities are present
  \end{cellist}
  &
  Successful deployment across all services; functionality verified across multiple edge cases. Tutor capability tested. Five instances of the same vulnerability generated to measure AI output variety and success rate.
  \\[4pt]

  10--12 &
  \begin{cellist}
    \item Analysis and report writing, evaluating results from the previous week (whether the AI consistently reproduced vulnerabilities)
  \end{cellist}
  &
  Final project report and source code submission ready.
  \\

\end{longtable}

\bib{bibliography}

\end{document}
